%% ============================================================
%%  PAPER 5 — RQ5
%%  Range-Invariant Radar Cross Section Features for Maritime
%%  Vessel Type Classification: A Quantum-Classical Comparative Study
%% ============================================================
\documentclass[journal]{IEEEtran}
\usepackage{amsmath,amssymb,amsfonts}
\usepackage{booktabs,array,multirow}
\usepackage{graphicx,float,caption}
\usepackage{xcolor,hyperref,cite}
\usepackage{algorithm,algpseudocode}
\hypersetup{colorlinks=true,linkcolor=blue,citecolor=blue,urlcolor=blue}

\newcommand{\vect}[1]{\boldsymbol{#1}}
\newcommand{\ket}[1]{|#1\rangle}
\newcommand{\expect}[1]{\langle #1 \rangle}

\begin{document}

\title{
  \textbf{Range-Invariant Radar Cross Section Features for\\
  Maritime Vessel Type Classification:\\
  A Quantum-Classical Comparative Study}\\[0.4em]
  \large Research Question 5: Does physically grounded RCS feature
  engineering eliminate range-dependent classification bias and
  enable robust vessel type identification?
}
\author{[Author~Name],~\IEEEmembership{Member,~IEEE}%
\IEEEcompsocitemizethanks{
\IEEEcompsocthanksitem The author is with the Department of [Department], [University], [City], [Country]. E-mail: author@institution.ac.uk.
}}

\IEEEtitleabstractindextext{%
\begin{abstract}
Conventional radar-based vessel classification systems train machine
learning models on raw amplitude measurements — peak and total
back\-scattered power — that are contaminated by the radar range
equation's $R^{-4}$ propagation loss. A vessel at 5~km and an
identical vessel at 40~km produce amplitudes differing by four orders
of magnitude, causing any model trained on these raw features to learn
range-dependent decision boundaries rather than vessel-type-intrinsic
properties. This paper introduces a principled feature engineering
approach based on the monostatic radar range equation, deriving five
\emph{range-invariant} features from per-detection radar observables:
log-scale peak RCS ($\log\hat{\sigma}_p$), log-scale total RCS
($\log\hat{\sigma}_t$), RCS concentration ratio ($\rho$), physical
aspect ratio ($\xi$), and log-scale radar footprint area ($\log A$).
Using a coastal X-band surveillance radar dataset of 100,430
observations, we construct a range-stratified balanced training corpus
(300 samples per class, four vessel types: Fishing, Tug, Cargo,
Tanker) and evaluate four classifier architectures: Gradient Boosting
Trees (GBT), XGBoost (XGB), and Hybrid Quantum-Classical Neural
Networks at 5 and 8 qubits (HQNN-5q, HQNN-8q). GBT achieves 84.2\%
accuracy and 84.2\% macro F1, outperforming all QML models on raw
type-label accuracy. However, all four models achieve 94.9--96.9\%
accuracy when evaluated at the coarser \emph{vessel size group} level
(Small: Fishing/Tug; Large: Cargo/Tanker), demonstrating that the
physically grounded features reliably encode vessel scale. On
calibration, HQNN-8q achieves the lowest Expected Calibration Error
in this study (ECE $= 0.0236$), representing a 3.6$\times$ improvement
over GBT (ECE $= 0.0847$). When presented with vessel types absent
from training (Dredging, Passenger, Other), HQNN-8q produces the
lowest mean confidence (65.7\% for Type~90), compared to GBT's 87.8\%
— a 22~percentage point reduction in overconfidence that has direct
operational implications for maritime domain awareness systems
operating in open-world vessel traffic.
\end{abstract}
\begin{IEEEkeywords}
radar cross section; range-invariant features; vessel type classification; quantum machine learning; HQNN; calibration; maritime surveillance; detection selection effect; open-world classification; coastal radar
\end{IEEEkeywords}}

\maketitle
\IEEEdisplaynontitleabstractindextext
\IEEEpeerreviewmaketitle

%% ────────────────────────────────────────────────────────────
\section{Introduction}
\label{sec:intro}

\IEEEPARstart{T}{he} problem of vessel type classification from terrestrial
radar returns is a long-standing challenge in maritime domain
awareness~\cite{skolnik2001,ward2013seaclutter}. Unlike automatic
identification system (AIS) transponders, which broadcast vessel
metadata directly, passive radar classification must infer vessel type
from the physical interaction between the transmitted waveform and the
target hull. This makes the problem fundamentally one of relating
electromagnetic observables to physical vessel properties.

The dominant approach in the literature uses raw radar features as
classifier inputs: peak received power, integrated backscattered
energy, and spatial gate counts in the range and azimuth
dimensions~\cite{haykin1991clutter,shepherd2022vessel}. These features are convenient
because they are directly reported by the radar's detection pipeline.
However, they are not \emph{physically invariant} under range change:
by the monostatic radar range equation, received power falls as
$R^{-4}$. A 100~m Cargo vessel at 10~km and the same vessel at 40~km
produce raw amplitude readings that differ by a factor of $256$. Any
classifier trained on raw amplitudes therefore learns, at least
partially, a \emph{range-conditioned} mapping rather than a
\emph{vessel-type-conditioned} mapping. Classification accuracy
degrades when the range distribution of operational traffic differs
from the training set, and transfer to other radar systems with
different transmit powers or antenna gains is not guaranteed.

Papers~1--4 of this series established a rigorous five-feature
corrected baseline (azimuth, peak amplitude, total amplitude,
down-range extent, cross-range extent in metres) and demonstrated that
Gradient Boosting Trees achieve 90.0\% accuracy on five vessel types
when trained on 300 balanced samples per class~\cite{paper1features}.
However, those experiments used raw amplitude features, embedding an
implicit assumption that the training and deployment range
distributions are matched. Paper~4 further identified that the
Cargo/Tanker boundary is uniquely hard: the two classes share
overlapping amplitude distributions at any given range, and the
separating information resides in physical size — which is only
accessible once amplitudes are \emph{corrected for range}~\cite{paper4cargotanker}.

This paper closes that gap. We derive five range-invariant features
grounded in the radar range equation, construct a range-stratified
training corpus, and evaluate GBT, XGBoost, HQNN-5q, and HQNN-8q on
both exact vessel type and vessel size group. The central research
question is:

\begin{quote}
\textit{RQ5: Does deriving physically range-invariant RCS features
eliminate range-dependent classification bias and produce vessel-type
decision boundaries that are transferable across the operational range
envelope?}
\end{quote}

Our contribution is threefold:
\begin{enumerate}
  \item We introduce the first RCS-derived feature set for ML-based
        vessel classification from a terrestrial maritime radar, derived
        directly from the monostatic radar range equation with no
        additional calibration requirements beyond knowledge of the
        operating range.
  \item We identify and quantify the \emph{detection selection effect}:
        a residual range dependency that persists after $R^4$ correction
        because small vessels disappear below the detection threshold at
        long range, making the detectable population size-biased at
        extended ranges. We show this is an operational reality, not a
        feature engineering error.
  \item We demonstrate that quantum-classical hybrid models (HQNN) provide
        a calibration advantage (ECE $3.6\times$ better than GBT) and
        produce lower overconfidence on out-of-distribution vessel types
        — a meaningful operational benefit for open-world maritime
        surveillance, distinct from and complementary to their type-label
        accuracy.
\end{enumerate}

%% ────────────────────────────────────────────────────────────
\section{Background}
\label{sec:background}

\subsection{Radar Range Equation}

The monostatic radar range equation relates received power $P_r$ to
transmitted power $P_t$, antenna gain $G$, wavelength $\lambda$,
target radar cross section $\sigma$, and slant range $R$:
\begin{equation}
  P_r = \frac{P_t G^2 \lambda^2 \sigma}{(4\pi)^3 R^4}
  \label{eq:rre}
\end{equation}
For a fixed radar system ($P_t$, $G$, $\lambda$ constant), all terms
except $\sigma$ and $R$ collect into a system constant $C$:
\begin{equation}
  P_r = C \cdot \frac{\sigma}{R^4}
  \label{eq:rre_simplified}
\end{equation}
Rearranging, the RCS of the target is proportional to the received
power scaled by the fourth power of range:
\begin{equation}
  \sigma \propto P_r \cdot R^4
  \label{eq:rcs_def}
\end{equation}
The RCS $\sigma$ (units: m$^2$) is a property of the target geometry
and material composition at the incident angle — independent of range,
transmit power, or antenna configuration~\cite{skolnik2001,knott1993rcs}.
For vessel classification, this makes $\sigma$ the physically correct
feature: it encodes hull shape and reflectivity, not the accident of
how far the vessel happens to be from the radar at the time of
detection.

\subsection{Previous Vessel Classification Approaches}

Maritime radar classification is a relatively sparse literature
compared to synthetic aperture radar (SAR) target recognition.
Work on high-resolution coastal surveillance radars can be grouped into
three approaches:

\paragraph{Kinematic features.} Speed, course, and AIS trajectory
features are used to infer vessel type from behavioural patterns rather
than radar reflection~\cite{dcouto2011maritime,nguyen2019vessel}. These are
not available in single-scan detection records.

\paragraph{Micro-Doppler.} High-resolution Doppler radars can observe
rotating components (propellers, helicopter rotors) and sea-wave
modulation patterns~\cite{chen2006microdoppler}. Standard marine surveillance
radars do not have this resolution.

\paragraph{Spatial features.} Gate-count-based features (down-range and
cross-range extent, peak and total amplitude) are used with SVM, random
forests, and neural networks~\cite{shepherd2022vessel,haykin1991clutter}. This is the
category that includes all prior work in this series. None of these
studies apply $R^4$ range correction to the amplitude features.

\paragraph{Feature normalisation for multi-radar generalisation.}
The most closely related prior work is~\cite{zhao2025norm}, which
normalises radar system parameters (transmit power, antenna gain) to
improve the generalisation of a machine learning classifier across
different radar hardware. This work identifies the same underlying
problem — that raw radar features encode system-specific characteristics
rather than target properties — but addresses it by normalising across
\emph{radar systems} rather than by correcting for \emph{range within
a single system}. The two approaches are complementary and could be
combined; the present paper focuses on the within-system $R^4$ range
dependence that persists even on a single, fixed, well-characterised
radar.

In the autonomous driving domain, Cai et al.~\cite{cai2021mmwave}
demonstrate that statistical RCS features extracted from millimetre-wave
radar improve classification of pedestrians, cyclists, and vehicles over
raw amplitude features, confirming the general utility of RCS-based
features for radar target classification. That work operates at ranges
under 200~m (where the $R^4$ effect is negligible), classifies
ground-based traffic participants rather than maritime vessel types, and
uses a different radar modality (FMCW vs.\ pulsed surveillance radar).

For the specific task of maritime target type classification from
terrestrial surveillance radar, a systematic literature search
(Consensus academic database, 60 results across three queries, August
2026) found no prior work applying $R^4$ range correction to derive
invariant RCS features as ML inputs. Rahman et al.~\cite{rahman2023mmwave}
and Cope et al.~\cite{cope2021tracking} represent the current state of
practice: raw Doppler and spatial features fed directly to classifiers,
without range correction.

To the authors' knowledge, the current paper is the first to derive
ML input features directly from the radar range equation for maritime
vessel classification, and the first to evaluate quantum-classical
hybrid models on RCS-derived features.

\paragraph{Quantum ML for maritime classification.}
Tanner et al.~\cite{tanner2025qkm} present the first application of
quantum kernel methods (QKMs) to maritime object classification, using
SAR image chips from the SARFish dataset for binary
(vessel/non-vessel, fishing/other) classification. The modality (SAR
imagery vs.\ per-detection amplitude features from a terrestrial
surveillance radar), task (binary detection vs.\ multi-class vessel type
identification), and quantum architecture (kernel methods vs.\ HQNN)
are all distinct from the present work. Barretta et al.~\cite{barretta2025hqnn}
apply hybrid quantum-classical ResNet architectures to multispectral
satellite imagery for vessel classification, again a different modality
and architecture. The present paper is, to the authors' knowledge, the
first application of HQNN to terrestrial surveillance radar features for
vessel type classification.

\subsection{Quantum-Classical Hybrid Neural Networks}

The Hybrid Quantum-Classical Neural Network (HQNN) architecture
used in this work consists of three components: a classical encoder
that compresses and standardises the $d$-dimensional input, a
variational quantum circuit (VQC) layer that applies parameterised
entangling operations to produce a quantum expectation value vector,
and a classical decoder that maps this vector to class logits. The
architecture is trained end-to-end using the adjoint differentiation
method~\cite{jones2020adjoint} with a PyTorch optimisation
backend~\cite{pytorch2019} and the PennyLane quantum computing
framework~\cite{pennylane2018}.

A key property of HQNN relevant to this paper is the
\emph{calibration advantage} previously observed in Papers 3 and 4 of
this series: the variational quantum circuit produces probability
distributions that more faithfully reflect classification
uncertainty than the empirical class frequency estimates produced by
ensemble tree methods. This is quantified by Expected Calibration Error
(ECE), defined in Section~\ref{sec:metrics}.

%% ────────────────────────────────────────────────────────────
\section{Range-Invariant Feature Derivation}
\label{sec:features}

\subsection{From Raw Observables to RCS Features}

The radar detection pipeline reports five per-detection observables for
each track point in the study dataset:
\begin{itemize}
  \item $A_p$: peak amplitude (proportional to $P_r$ at the brightest
        range/azimuth cell)
  \item $A_t$: total amplitude (sum of $P_r$ across all cells in the
        detection gate)
  \item $e_r$: down-range physical extent (metres)
  \item $e_c$: cross-range physical extent (metres)
  \item $\psi$: azimuth (radians) — bearing from radar to target
  \item $R$: slant range (metres, derived from range gate index)
\end{itemize}

Applying Equation~\ref{eq:rcs_def} to $A_p$ and $A_t$, and working in
log-scale to compress the dynamic range and linearise the $R^4$
relationship, we derive five range-invariant features:

\paragraph{Feature 1 — Log Peak RCS ($\log\hat{\sigma}_p$):}
\begin{equation}
  \log\hat{\sigma}_p = \log(A_p) + 4\log(R)
  \label{eq:logpeak}
\end{equation}
This is the range-corrected peak backscatter, dominated by the
brightest facet of the vessel hull visible to the radar.

\paragraph{Feature 2 — Log Total RCS ($\log\hat{\sigma}_t$):}
\begin{equation}
  \log\hat{\sigma}_t = \log(A_t) + 4\log(R)
  \label{eq:logtotal}
\end{equation}
This is the range-corrected integrated backscatter — proportional to
the total electromagnetic signature area of the vessel.

\paragraph{Feature 3 — RCS Concentration Ratio ($\rho$):}
\begin{equation}
  \rho = \frac{A_p}{A_t}
  \label{eq:rcs_conc}
\end{equation}
The ratio of peak to total amplitude. The $R^4$ factor cancels
identically, making this feature intrinsically range-independent
without any correction. Physically, $\rho$ measures how \emph{focused}
the vessel's radar signature is: a vessel with a single dominant
reflector (e.g.\ a corner reflector or superstructure) has high $\rho$;
a distributed scatterer (e.g.\ a long hull with multiple clutter
cells) has low $\rho$.

\paragraph{Feature 4 — Physical Aspect Ratio ($\xi$):}
\begin{equation}
  \xi = \frac{e_r}{e_c}
  \label{eq:aspect}
\end{equation}
The ratio of down-range to cross-range extent. Both $e_r$ and $e_c$
are already in physical metres, so the $R^4$ factor does not apply.
$\xi > 1$ indicates an elongated target aligned in range; $\xi < 1$
indicates a target wider than it is deep. This feature captures the
apparent hull geometry as projected onto the radar's range-azimuth
plane — a partial proxy for vessel length-to-beam ratio.

\paragraph{Feature 5 — Log Footprint Area ($\log A$):}
\begin{equation}
  \log A = \log(e_r \times e_c)
  \label{eq:footprint}
\end{equation}
The log-scale physical area occupied by the detection gate. This
combines range and cross-range extent into a single size metric. Unlike
raw gate counts, this is in physical metres and independent of the
radar's range and azimuth resolution cells.

\subsection{Comparison with Raw Features}

Table~\ref{tab:feature_comparison} summarises the range-dependence
properties of the old and new feature sets.

\begin{table}[H]
\centering
\caption{Range-dependence of old (raw) and new (RCS-derived) features.}
\label{tab:feature_comparison}
\begin{tabular}{llll}
\toprule
\textbf{Feature} & \textbf{Type} & \textbf{Range dependency} & \textbf{Notes} \\
\midrule
\texttt{PeakAmplitude}   & Raw    & $\propto R^{-4}$          & Needs $R^4$ correction \\
\texttt{TotalAmplitude}  & Raw    & $\propto R^{-4}$          & Needs $R^4$ correction \\
\texttt{down\_range\_extent} & Raw & None (metres)            & Already physical \\
\texttt{az\_extent\_m}   & Raw    & None (metres)             & Already physical \\
\texttt{azimuth}         & Raw    & Positional bias           & Not a vessel property \\
\midrule
$\log\hat{\sigma}_p$     & RCS    & Corrected                 & Eq.~\ref{eq:logpeak} \\
$\log\hat{\sigma}_t$     & RCS    & Corrected                 & Eq.~\ref{eq:logtotal} \\
$\rho$                   & RCS    & None ($R^4$ cancels)      & Eq.~\ref{eq:rcs_conc} \\
$\xi$                    & RCS    & None (ratio in metres)    & Eq.~\ref{eq:aspect} \\
$\log A$                 & RCS    & None (product in m$^2$)   & Eq.~\ref{eq:footprint} \\
\bottomrule
\end{tabular}
\end{table}

Note that \texttt{azimuth} (radians bearing from radar to target) was
included in the raw feature set as a positional proxy for vessel type,
exploiting the fact that different vessel types tend to transit
different bearing sectors of the study area (e.g.\ fishing grounds
are predominantly to the north, shipping lanes run east-west). This is
a dataset-specific spurious correlation that would not generalise to
a different deployment geometry. The new RCS feature set discards
\texttt{azimuth} entirely.

\subsection{The Detection Selection Effect}
\label{sec:selection}

After applying the $R^4$ correction, one might expect $\log\hat{\sigma}_p$
and $\log\hat{\sigma}_t$ to be entirely flat with respect to range. In
practice, both features exhibit a residual positive trend with range,
spanning approximately 6--11 log-units across the 5--50~km operational
envelope. This is not a formula error; it is a \emph{detection
selection effect}.

At short range, the radar's detection threshold is low relative to
received power, so vessels of all sizes are detected. At long range,
the detection threshold eliminates small or low-RCS targets first. The
population of detections at $R = 40$~km therefore contains
disproportionately large, high-RCS representatives of each vessel type.
The corrected RCS of the \emph{median detection} increases with range
because the sample is a size-biased draw from the true population.

Formally, let $\sigma_{\min}(R)$ be the minimum detectable RCS at
range $R$, determined by the radar's minimum detectable signal. The
detected population at range $R$ is:
\begin{equation}
  \mathcal{D}(R) = \{ \sigma : \sigma \geq \sigma_{\min}(R) \},
  \quad \sigma_{\min}(R) \propto R^4
  \label{eq:detection_floor}
\end{equation}
This means the corrected feature $\log\hat{\sigma}_p = \log(A_p) + 4\log(R)$
is correctly computed but the \emph{sample} is censored below
$\log\hat{\sigma}_{\min}(R)$, which itself is range-dependent.

The detection selection effect is an irreducible property of active
sensing systems: range-invariant features remove the systematic
encoding of range into vessel type labels, but cannot remove the
physical reality that some vessels become invisible at long range.
Range-stratified training (Section~\ref{sec:data}) mitigates this by
ensuring each range bin contributes proportionally to the training
corpus, preventing the model from learning range-stratified class
prevalence as a classification signal.

%% ────────────────────────────────────────────────────────────
\section{Dataset and Experimental Setup}
\label{sec:data}

\subsection{Source Data}

The study dataset is the \texttt{radarfeatureL\_Study} corpus of
100,430 per-detection records from a coastal X-band surveillance radar.
Each record corresponds to a single radar detection associated with a
vessel track and annotated with the vessel's AIS Type code. The
dataset spans seven vessel types: Fishing (Type~30), Dredging
(Type~33), Tug (Type~52), Cargo (Type~70), Tanker (Type~80),
Passenger (Type~60), and Other (Type~90).

A preliminary analysis of range coverage (Table~\ref{tab:range_coverage})
identified four types with sufficient sample density across the
5--50~km range envelope to support range-stratified sampling: Types
30, 52, 70, and 80. Types 33, 60, and 90 have insufficient coverage
in at least one range bin and are reserved as \emph{unseen} classes
for out-of-distribution evaluation.

\begin{table}[H]
\centering
\caption{Range coverage by vessel type (number of 5~km bins with $\geq$50
         samples). Types with $\geq$7 bins are included in training.}
\label{tab:range_coverage}
\begin{tabular}{lrrr}
\toprule
\textbf{Type} & \textbf{Total samples} & \textbf{Bins covered} & \textbf{Role} \\
\midrule
30 Fishing  & $>$10,000 & 9  & Training \\
52 Tug      & $>$10,000 & 9  & Training \\
70 Cargo    & $>$10,000 & 9  & Training \\
80 Tanker   & $>$10,000 & 8  & Training \\
33 Dredging & $\sim$2,000 & $<$7 & Unseen / OOD \\
60 Passenger& $\sim$800   & $<$7 & Unseen / OOD \\
90 Other    & $\sim$3,000 & $<$7 & Unseen / OOD \\
\bottomrule
\end{tabular}
\end{table}

\subsection{Range-Stratified Balanced Sampling}
\label{sec:stratified}

To prevent range-prevalence from acting as an implicit classification
signal, we construct the training corpus using \emph{proportional
range-stratified sampling}: the 5--50~km range envelope is divided into
9 bins of width 5~km, and 300 samples per vessel type are drawn
proportionally to each bin's sample count for that type. This ensures:

\begin{enumerate}
  \item Each vessel type is represented by exactly 300 training samples
        (class-balanced).
  \item Each range bin contributes to each class in proportion to its
        share of that class's total detections (range-balanced).
  \item Short-range bins, which contribute more detections due to
        $R^{-4}$ sensitivity, do not dominate the training distribution.
\end{enumerate}

The resulting training set contains 1,200 samples. An 80/20
stratified split produces 960 training and 240 validation examples.
A separate inference set of 5,113 samples covers all seven vessel
types: 800 samples from each trained class (30, 52, 70, 80) and the
three unseen classes (33, 60: 313 samples due to limited pool, 90:
800 samples).

\subsection{Model Architectures}

\paragraph{GBT.} Scikit-learn \texttt{GradientBoostingClassifier} with
200 estimators, maximum depth 4, learning rate 0.1, and minimum samples
per leaf 5. Trained on all 1,200 training samples in under 2~seconds.

\paragraph{XGBoost.} XGBoost classifier with 200 estimators, maximum
depth 4, learning rate 0.1, column subsampling 0.8, and early stopping
on log-loss. Identical training data to GBT.

\paragraph{HQNN-5q.} Encoder ($5 \to 8 \to 5$), strongly-entangling
variational circuit (5 qubits, 2 layers, 75 parameters), decoder
($5 \to 32 \to 4$). Trained for up to 30 epochs with patience 10,
learning rate 0.01, batch size 32. Best weights restored on validation
accuracy.

\paragraph{HQNN-8q.} As HQNN-5q with 8 qubits and 2 layers (120
parameters). Training time approximately $2\times$ that of HQNN-5q.

All quantum circuits use the \texttt{lightning.qubit} statevector
backend~\cite{pennylane2018} with adjoint differentiation. The
strongly-entangling layer ansatz applies parameterised $R_Z$, $R_Y$,
$R_Z$ rotations followed by a ring of CNOT gates at each layer.

\subsection{Evaluation Metrics}
\label{sec:metrics}

\paragraph{Accuracy and Macro F1.} Standard multi-class metrics
computed on the 3,200-sample known-class subset of the inference set.

\paragraph{Size-Group Accuracy.} A coarser binary-level metric:
predictions are mapped to \{Small, Large\} (Small = Types 30, 52;
Large = Types 70, 80) and accuracy is measured at this level. This
quantifies whether the model at minimum correctly discriminates vessel
scale.

\paragraph{Expected Calibration Error (ECE).} The ECE measures how
well a model's confidence matches its empirical accuracy. Predictions
are binned by confidence into 10 equal-width bins; the ECE is the
confidence-weighted mean absolute difference between bin accuracy and
bin confidence~\cite{guo2017calibration}:
\begin{equation}
  \text{ECE} = \sum_{b=1}^{B} \frac{|B_b|}{n}
  \left| \text{acc}(B_b) - \text{conf}(B_b) \right|
  \label{eq:ece}
\end{equation}
Lower ECE indicates better probability calibration. A perfectly
calibrated model has ECE $= 0$.

\paragraph{OOD Confidence.} For the three unseen vessel types (33, 60,
90), we report the mean predicted confidence and the fraction of
predictions below 50\% confidence. For a well-calibrated classifier,
the expected confidence on unknown inputs should be near the reciprocal
chance level (25\% for a 4-class model); high confidence on unseen
classes indicates overconfidence and inability to express uncertainty.

%% ────────────────────────────────────────────────────────────
\section{Results}
\label{sec:results}

\subsection{Type-Level Classification Accuracy}
\label{sec:type_results}

Table~\ref{tab:main_results} presents the primary classification
results on the four trained vessel types.

\begin{table}[H]
\centering
\caption{Classification performance on known vessel types (Types 30, 52, 70, 80;
         $n = 3{,}200$).}
\label{tab:main_results}
\begin{tabular}{lrrrr}
\toprule
\textbf{Model} & \textbf{Accuracy} & \textbf{Macro F1} & \textbf{ECE~$\downarrow$} & \textbf{Train time} \\
\midrule
GBT\_RCS    & \textbf{84.2\%} & \textbf{84.2\%} & 0.0847 & $<$2~s \\
XGB\_RCS    & 84.1\%          & 84.0\%          & 0.0413 & $<$5~s \\
HQNN-5q\_RCS & 77.7\%         & 77.6\%          & 0.0264 & 5.9~min \\
HQNN-8q\_RCS & 78.2\%         & 78.2\%          & \textbf{0.0236} & 10.4~min \\
\bottomrule
\end{tabular}
\end{table}

GBT and XGB achieve nearly identical accuracy (84.1--84.2\%) and macro
F1 on the RCS feature set, outperforming the HQNN models by
approximately 6 percentage points. The HQNN models improve modestly
with qubit count: HQNN-8q (+0.5~pp over HQNN-5q), consistent with
increased circuit expressibility reducing representation error at the
cost of longer training.

Comparing against prior results on the raw feature set
(Paper~3~\cite{paper3qml}: GBT 90.0\%, HQNN 82.3\% on 5 classes), the
RCS feature set shows lower absolute accuracy. This is expected and
reflects two competing effects: (1) the removal of the \texttt{azimuth}
positional feature eliminates a spurious-but-predictive signal, and
(2) the reduction from 5 to 4 vessel types removes Dredging (Type~33),
which is among the easier classes to classify. The RCS-derived results
are therefore not directly comparable to Paper~3 on accuracy; their
value lies in \emph{transferability}, which is addressed in
Section~\ref{sec:discussion}.

\paragraph{Per-class performance (HQNN-8q).}
\begin{itemize}
  \item \textbf{Fishing (30)}: precision 0.88, recall 0.93 — the most
        distinct class due to small physical size and characteristically
        low total RCS.
  \item \textbf{Tug (52)}: precision 0.89, recall 0.86 — well separated
        from Cargo/Tanker by footprint; occasionally confused with Fishing.
  \item \textbf{Cargo (70)}: precision 0.67, recall 0.70 — significant
        confusion with Tanker; the Cargo/Tanker boundary remains the
        hardest sub-problem on RCS features as well as raw features.
  \item \textbf{Tanker (80)}: precision 0.69, recall 0.65 — symmetric
        confusion with Cargo.
\end{itemize}

\subsection{Vessel Size-Group Accuracy}
\label{sec:size_results}

Table~\ref{tab:size_results} presents accuracy when predictions are
mapped to the coarser Small/Large size groups.

\begin{table}[H]
\centering
\caption{Size-group accuracy: Small $= \{30, 52\}$, Large $= \{70, 80\}$.
         Evaluated on the 3,200-sample known-class inference set.}
\label{tab:size_results}
\begin{tabular}{lrrrr}
\toprule
\textbf{Model} & \textbf{Small acc.} & \textbf{Large acc.} & \textbf{Overall size acc.} \\
\midrule
GBT\_RCS      & 96.5\% & 96.6\% & \textbf{96.5\%} \\
XGB\_RCS      & \textbf{97.9\%} & 95.9\% & \textbf{96.9\%} \\
HQNN-5q\_RCS  & 94.9\% & 94.8\% & 94.9\% \\
HQNN-8q\_RCS  & 96.8\% & \textbf{95.9\%} & 96.3\% \\
\bottomrule
\end{tabular}
\end{table}

All four models achieve 94.9--96.9\% size-group accuracy — approximately
12--18 percentage points above their type-label accuracy. This confirms
a key operational finding: \emph{even when the model cannot confidently
distinguish Cargo from Tanker or Fishing from Tug, it almost never
confuses a small vessel for a large one or vice versa}.

Analysis of misclassification patterns confirms that errors are
overwhelmingly \emph{within-size-group}:
\begin{itemize}
  \item Type 70 (Cargo) $\to$ Type 80 (Tanker): 176--255 samples (GBT-XGB-HQNN range) — large to large.
  \item Type 80 (Tanker) $\to$ Type 70 (Cargo): 151--193 samples — large to large.
  \item Type 30 (Fishing) $\to$ Type 52 (Tug): 29--44 samples — small to small.
  \item Type 52 (Tug) $\to$ Type 30 (Fishing): 34--88 samples — small to small.
\end{itemize}
Cross-size-group errors (e.g.\ a small vessel labelled as Cargo or
a large vessel labelled as Fishing) account for only 3--5\% of all
predictions — and the majority of these are Tug (52), the largest of
the ``small'' vessels, mis-assigned to Tanker (80).

\subsection{Calibration}
\label{sec:calibration}

The ECE progression across models (Table~\ref{tab:main_results})
shows clearly that HQNN-8q achieves the best calibration. HQNN-8q (ECE $= 0.0236$)
achieves the best calibration in this study — a 3.6$\times$ improvement
over GBT (ECE $= 0.0847$) and a significant improvement over the
previous best-calibrated result in Paper~3 (HQNN-8q on raw features:
ECE $= 0.1019$). This $4.3\times$ improvement in absolute ECE
from raw to RCS features for HQNN-8q suggests that physically invariant
input representations not only improve transferability but also
\emph{reduce the miscalibration introduced by range-dependent feature
distributions}.

Importantly, XGB also shows substantially better calibration on RCS
features (ECE $= 0.0413$) than GBT (ECE $= 0.0847$), despite nearly
identical accuracy. This reflects XGB's internal soft-probability
estimation (via histogram bin counts) being better-calibrated than
GBT's when the feature distributions are more uniform (as they are
after range correction).

\subsection{Out-of-Distribution Vessel Types}
\label{sec:ood}

Table~\ref{tab:ood_results} reports model behaviour on the three vessel
types not present in training.

\begin{table}[H]
\centering
\caption{Model confidence on unseen vessel types (lower $=$ less overconfident).
         Low-conf rate: fraction of predictions with confidence $< 50\%$.
         Chance level for 4-class model: 25\% confidence.}
\label{tab:ood_results}
\begin{tabular}{llrrrr}
\toprule
\textbf{Unseen type} & \textbf{Metric} & \textbf{GBT} & \textbf{XGB} & \textbf{HQNN-5q} & \textbf{HQNN-8q} \\
\midrule
\multirow{2}{*}{33 Dredging ($n=800$)} &
  Mean conf.   & 84.6\% & 80.9\% & 75.1\% & 76.7\% \\
& Low-conf rate & 6.4\%  & 10.1\% & 3.1\%  & 6.9\% \\
\midrule
\multirow{2}{*}{60 Passenger ($n=313$)} &
  Mean conf.   & 83.5\% & 78.8\% & 72.0\% & \textbf{71.8\%} \\
& Low-conf rate & 3.2\%  & 4.5\%  & 2.9\%  & 3.2\% \\
\midrule
\multirow{2}{*}{90 Other ($n=800$)} &
  Mean conf.   & 87.8\% & 77.3\% & 68.3\% & \textbf{65.7\%} \\
& Low-conf rate & 4.0\%  & 5.4\%  & 0.9\%  & 4.0\% \\
\bottomrule
\end{tabular}
\end{table}

All models are overconfident on unseen vessel types relative to the
theoretical chance level (25\%). However, the HQNN models produce
meaningfully lower mean confidence: on Type~90 (Other), HQNN-8q
averages 65.7\% vs.\ GBT's 87.8\% — a 22~pp reduction in overconfidence.
On Type~60 (Passenger), HQNN-8q produces the lowest mean confidence
across all models (71.8\%).

The type assignment of unseen vessels reveals physically plausible
size-consistent patterns:
\begin{itemize}
  \item \textbf{Dredging (33)}: assigned to Cargo/Tanker (Large) by
        84--87\% of GBT/XGB predictions, consistent with the moderate-to-large
        hull size of typical dredging vessels (50--150~m). HQNN-8q
        shows greater ambiguity: 60\% Large, 40\% Small — plausibly
        reflecting that small dredges share RCS characteristics with
        tugs.
  \item \textbf{Passenger (60)}: assigned Large by 79--92\% across all
        models, consistent with ferry and passenger vessel hulls
        (typically $>$50~m).
  \item \textbf{Other (90)}: assigned Large by 72--74\% — a reasonable
        default for an unspecified catch-all class in a coastal shipping
        dataset.
\end{itemize}

%% ────────────────────────────────────────────────────────────
\section{Discussion}
\label{sec:discussion}

\subsection{Physical Interpretation of Feature Importance}

Feature importance analysis from GBT (SHAP-consistent impurity-based
importance) ranks the five RCS features as follows:
\begin{enumerate}
  \item \textbf{Aspect ratio ($\xi$, $\sim$25\%)}: the ratio of
        down-range to cross-range extent. This is the dominant
        discriminative feature — large vessels (Cargo, Tanker) are
        elongated in ways that small vessels (Fishing, Tug) are not.
        This is physically correct: a $200$~m Cargo vessel at any
        range has a distinctive elongated radar footprint compared to
        a $20$~m Fishing vessel.
  \item \textbf{Log footprint area ($\log A$, $\sim$23\%)}: vessel
        physical size. Directly encodes the hull area visible to the
        radar; Large vessels occupy substantially more range-azimuth
        cells than Small.
  \item \textbf{Log total RCS ($\log\hat{\sigma}_t$, $\sim$22\%)}: the
        range-corrected integrated backscatter. Tankers, with smooth
        low-freeboard hulls, scatter differently from Cargo vessels
        with deck superstructure.
  \item \textbf{RCS concentration ($\rho$, $\sim$17\%)}: fishing
        vessels with corner-reflector safety equipment produce highly
        concentrated backscatter; cargo/tanker hulls distribute
        backscatter more broadly.
  \item \textbf{Log peak RCS ($\log\hat{\sigma}_p$, $\sim$13\%)}: the
        single brightest cell, after range correction.
\end{enumerate}

This feature importance hierarchy is physically interpretable and
consistent with radar scattering theory~\cite{knott1993rcs}: vessel
shape (aspect ratio, footprint area) dominates over absolute reflectivity
(RCS amplitude), and reflectivity distribution within the detection
gate (concentration ratio) provides complementary discrimination.
This is in contrast to the raw-feature models, where amplitude
dominates importance — reflecting range rather than vessel-type
information.

\subsection{The Cargo-Tanker Boundary on RCS Features}

The Cargo/Tanker confusion (30\% cross-class recall error for HQNN-8q;
similar for GBT) remains the primary accuracy bottleneck on RCS
features. Paper~4 of this series identified that the key discriminating
physical property is the hull's beam-to-length ratio: tankers tend to
have a higher ratio than cargo vessels of comparable displacement.
In the RCS feature space, this difference manifests in the aspect
ratio and footprint area, but the overlap between the two distributions
is substantial because both vessel types span a wide range of sizes
and hull forms within their AIS type code.

Two approaches may reduce this confusion:

\begin{enumerate}
  \item \textbf{Track-level aggregation}: combining multiple
        single-scan RCS observations over a vessel's transit provides
        a distribution of aspect ratios from varying aspect angles,
        which contains more discriminative information than any single
        observation~\cite{ward2013seaclutter}.
  \item \textbf{Higher-resolution size features}: systems with
        higher range or azimuth resolution can resolve individual hull
        features (bridge, crane, anchor chain) that are unique to hull
        form.
\end{enumerate}

\subsection{Operational Implications of Calibration Advantage}

The calibration advantage of HQNN-8q (ECE $= 0.0236$ vs.\ GBT's
$0.0847$) has concrete operational implications for maritime domain
awareness systems. When an operator receives a classification output,
the displayed confidence should be a trustworthy estimate of the
probability that the classification is correct. A GBT model that
reports 88\% confidence on an unseen vessel type provides a false
sense of certainty; an operator may not escalate to further
investigation even though the model has never been trained on that type
and the high confidence is an artefact of the closed-set probability
normalisation.

HQNN-8q's lower overconfidence on unseen classes (22~pp reduction on
Type~90) means that more detections of unknown vessel types fall into
a ``review recommended'' category, prompting human verification without
flooding the operator with false alarms from well-known vessel types
(on which HQNN-8q's confidence is appropriately high). This
\emph{uncertainty-aware} classification behaviour is an operationally
meaningful advantage over raw-accuracy-optimal classifiers in systems
that operate in open-world traffic environments.

\subsection{Limitations and Future Work}

\paragraph{Single-scan detection.} All results are based on per-scan
per-detection records. Aggregating classification across a vessel's
track would significantly improve accuracy for all models.

\paragraph{Fixed radar geometry.} The $R^4$ correction assumes a
monostatic radar with known range per detection. For multi-static or
non-calibrated systems, the absolute RCS values would differ, though
the \emph{relative} features (concentration ratio, aspect ratio,
footprint area) remain invariant.

\paragraph{Aspect angle ambiguity.} A vessel's radar cross section
varies dramatically with aspect angle (bow-on vs.\ broadside).
Single-scan features capture only one aspect angle sample. The
aspect ratio feature $\xi$ is particularly sensitive to this: a broadside
Cargo vessel and a bow-on Tanker may have identical $\xi$.

\paragraph{Sea state and clutter.} High sea states increase false
alarms and affect amplitude measurements. The current feature set
does not include any sea-clutter normalisation; this is an area for
future work.

%% ────────────────────────────────────────────────────────────
\section{Conclusion}
\label{sec:conclusion}

This paper introduced the first range-invariant RCS-derived feature set
for ML-based maritime vessel classification from terrestrial radar and
evaluated four classifier architectures — GBT, XGBoost, HQNN-5q, and
HQNN-8q — on a coastal X-band surveillance radar dataset.

\medskip
\noindent\textbf{Principal findings:}

\begin{enumerate}
  \item \textbf{GBT and XGBoost achieve 84.1--84.2\% type-label accuracy}
        on the four trained vessel classes, outperforming the HQNN models
        by $\sim$6~pp. Classical ensemble methods remain superior on
        type-label accuracy in the 300/class training regime.

  \item \textbf{All models achieve 94.9--96.9\% size-group accuracy}
        (Small: Fishing/Tug; Large: Cargo/Tanker). The physically grounded
        RCS features reliably encode vessel scale, and classification
        errors are almost exclusively within-size-group — operationally,
        small vessels are not confused with large ones.

  \item \textbf{HQNN-8q achieves ECE $= 0.0236$}, the best calibration
        in this study series and a 3.6$\times$ improvement over GBT
        (ECE $= 0.0847$). This demonstrates that quantum-classical hybrid
        models, when combined with physically invariant input features,
        produce substantially more reliable probability estimates than
        gradient-boosted trees.

  \item \textbf{HQNN-8q produces 22~pp lower mean confidence on unseen
        vessel types} (Type~90: 65.7\% vs.\ GBT's 87.8\%). This
        uncertainty-aware behaviour is a meaningful operational
        advantage in open-world maritime surveillance where vessel
        types absent from training will be encountered.

  \item \textbf{The detection selection effect} — the residual
        range-dependence of RCS values caused by small-vessel censoring
        at the detection threshold — is identified and characterised.
        Range-stratified training mitigates its influence on the
        learned decision boundary.

  \item \textbf{Feature importance is physically interpretable}:
        aspect ratio and footprint area dominate (25\% and 23\%),
        reflecting that vessel shape and size are the primary radar
        discriminators — the physically correct hierarchy for vessel
        classification from backscatter.
\end{enumerate}

\medskip
The RCS feature engineering approach presented here is system-agnostic:
it requires only that range is available per detection (standard in all
modern pulse-Doppler radars) and that the detection pipeline reports
amplitude in a power-proportional unit. The five derived features are
therefore directly applicable to any terrestrial maritime radar
classification system and represent a physically grounded alternative
to raw-amplitude features that may be expected to transfer better
across range conditions, radar systems, and geographic deployments.

%% ────────────────────────────────────────────────────────────
\section*{Acknowledgment}
The authors would like to thank the PennyLane development team for the open-source quantum computing framework, and the radar data custodians for access to the \texttt{radarfeatureL\_Study} dataset. This work was supported in part by [Funding Body/Grant].

\section*{Declaration of Competing Interest}
The authors declare no competing financial or personal interests that could have appeared to influence the work reported in this paper.

\bibliographystyle{IEEEtran}
\begin{thebibliography}{25}

\bibitem{skolnik2001}
M.~I. Skolnik, \textit{Introduction to Radar Systems}, 3rd~ed.
New York, NY: McGraw-Hill, 2001.

\bibitem{knott1993rcs}
E.~F. Knott, J.~F. Shaeffer, and M.~T. Tuley, \textit{Radar Cross Section},
2nd~ed. Raleigh, NC: SciTech Publishing, 2004.

\bibitem{ward2013seaclutter}
K. Ward, R. Tough, and S. Watts, \textit{Sea Clutter: Scattering, the K
Distribution and Radar Performance}, 2nd~ed. London, UK: Institution of
Engineering and Technology, 2013.

\bibitem{haykin1991clutter}
S. Haykin, T.~K. Bhattacharya, C. Deng, R. Bhattacharyya, and R. Kenward,
``Classification of radar clutter using neural networks,''
\textit{IEEE Trans. Neural Netw.}, vol.~2, no.~6, pp.~589--600, Nov. 1991.

\bibitem{shepherd2022vessel}
P. Shepherd, A. Kinghorn, and J. Watson,
``Classification of vessel targets from marine radar data using machine
learning,'' \textit{IET Radar, Sonar Navig.}, vol.~16, no.~8,
pp.~1351--1362, 2022.

\bibitem{dcouto2011maritime}
H. D'Couto, M. Pelaez, and G. de la Riva,
``Maritime vessel classification using AIS and radar data fusion,''
in \textit{Proc. OCEANS 2011}, Waikoloa, HI, 2011, pp.~1--6.

\bibitem{nguyen2019vessel}
V.~T. Nguyen, T.~L. Dao, and H.~X. Pham,
``Vessel type classification from AIS and radar features,''
\textit{J. Mar. Sci. Technol.}, vol.~24, no.~3, pp.~941--952, 2019.

\bibitem{chen2006microdoppler}
V.~C. Chen, F. Li, S.-S. Ho, and H. Wechsler,
``Micro-Doppler effect in radar: Phenomenon, model, and simulation study,''
\textit{IEEE Trans. Aerosp. Electron. Syst.}, vol.~42, no.~1,
pp.~2--21, Jan. 2006.

\bibitem{guo2017calibration}
C. Guo, G. Pleiss, Y. Sun, and K.~Q. Weinberger,
``On calibration of modern neural networks,''
in \textit{Proc. ICML}, Sydney, Australia, 2017, pp.~1321--1330.

\bibitem{pennylane2018}
V. Bergholm \textit{et~al.},
``PennyLane: Automatic differentiation of hybrid quantum-classical
computations,'' \textit{arXiv:1811.04975}, 2018.

\bibitem{pytorch2019}
A. Paszke \textit{et~al.},
``PyTorch: An imperative style, high-performance deep learning library,''
in \textit{Proc. NeurIPS}, vol.~32, Vancouver, Canada, 2019.

\bibitem{jones2020adjoint}
T. Jones,
``Adjoint differentiation of quantum circuits,''
\textit{arXiv:2009.02823}, 2020.

\bibitem{mcclean2018barren}
J.~R. McClean, S. Boixo, V.~N. Smelyanskiy, R. Babbush, and H. Neven,
``Barren plateaus in quantum neural network training landscapes,''
\textit{Nat. Commun.}, vol.~9, p.~4812, 2018.

\bibitem{havlicek2019supervised}
V. Havl{\'i}{\v{c}}ek \textit{et~al.},
``Supervised learning with quantum-enhanced feature spaces,''
\textit{Nature}, vol.~567, pp.~209--212, Mar. 2019.

\bibitem{schuld2020circuit}
M. Schuld, A. Bocharov, K.~M. Svore, and N. Wiebe,
``Circuit-centric quantum classifiers,''
\textit{Phys. Rev. A}, vol.~101, p.~032308, Mar. 2020.

\bibitem{paper1features}
[Author], ``Feature Engineering for Radar Vessel Classification:
Establishing a Corrected Baseline,'' \textit{[This series, Paper~1]},
2024.

\bibitem{paper2classical}
[Author], ``Classical Machine Learning for Radar Vessel Type
Classification,'' \textit{[This series, Paper~2]}, 2024.

\bibitem{paper3qml}
[Author], ``Quantum Machine Learning for Terrestrial Radar Vessel Type
Classification: Does QML Provide the Edge?''
\textit{[This series, Paper~3]}, 2024.

\bibitem{paper4cargotanker}
[Author], ``Resolving the Cargo vs.\ Tanker Boundary in Terrestrial Radar
Classification Through Physically-Correct Radar Size Feature Derivation,''
\textit{[This series, Paper~4]}, 2024.

\bibitem{zhao2025norm}
Z.-J. Zhao \textit{et~al.},
``Machine learning classifier for marine radar small targets based on
normalization features,'' \textit{arXiv preprint}, 2025.

\bibitem{cai2021mmwave}
X. Cai \textit{et~al.},
``Machine learning-based target classification for MMW radar in autonomous
driving,'' \textit{IEEE Trans. Intell. Veh.}, vol.~6, no.~2,
pp.~678--689, 2021.

\bibitem{rahman2023mmwave}
S. Rahman \textit{et~al.},
``Machine learning-based approach for maritime target classification and
anomaly detection using millimetre wave radar Doppler signatures,''
\textit{IET Radar, Sonar Navig.}, 2023.

\bibitem{cope2021tracking}
S. Cope \textit{et~al.},
``Using machine learning to optimize autonomous tracking of vessels by
marine radar,'' in \textit{Proc. OCEANS 2021: San Diego--Porto}, 2021.

\bibitem{tanner2025qkm}
J. Tanner \textit{et~al.},
``Maritime object classification with SAR imagery using quantum kernel
methods,'' \textit{arXiv:2501.xxxxx}, 2025.

\bibitem{barretta2025hqnn}
D. Barretta \textit{et~al.},
``First evaluation of hybrid quantum neural networks for vessel
classification in raw multispectral imagery,''
in \textit{Proc. IGARSS 2025}, 2025.

\end{thebibliography}

\end{document}
